% ch4_a §4.1–4.2 (书页 119–128 / p133–p142)
% 第4章 STATIC PROPERTIES OF SOLIDS

\chapter{固体的静态性质}

\epigraphbook{停住吧，你们这永远运转的天球……}{MARLOWE}

\section{固体的类型：能带图像}

我们已经花了相当长的篇幅讨论电子结构理论，因为它决定了固体的类型及其宏观性质。

考虑一条完全填满的能带，其上方存在一个能隙。这就是基态；由对称性可知，它必然处于这样的状态：没有电流在流动。要产生电流，可以施加一个电场，它把一些电子激发到代表净电流的态上。但如果存在能隙，就需要有限的激发能才能把电子抬升、越过能隙进入上面一条能带。这份能量无法由小的恒定电场提供，所以这种固体将是\emph{绝缘体}（insulator）。

%FIG{67}: {图 67　(a) 半导体中载流子的激发。(b) 金属中传输电流的电子。}

但假定能隙 $\mathcal{E}_{\mathrm{gap}}$ 很小。在温度 $T$ 下，将会有一个小而非零的电子密度，它们被热涨落激发到上面的能带中去——这一密度由形如 $\exp(-\mathcal{E}_{\mathrm{gap}}/kT)$ 的 Boltzmann 因子刻画。这些电子很容易传输电流，于是材料将具有可观察的电导率，并且随温度升高迅速增大。这时我们就说，我们得到的是\emph{半导体}（semiconductor）。

如果一条能带并未填满，例如在一价金属中，那么尽管系统的真实基态在 $\psi_{\mathbf{k}}$ 与 $\psi_{-\mathbf{k}}$ 之间是对称的、因而同样不载流，但在占据能级顶端之上能量距离为无穷小处，就存在着载流的态。此时电导率很高，而且对温度的依赖不大——温度只是在它支配电子散射机制这一点上才起作用。这样我们便得到了典型的\emph{金属}。

要判断一个具有给定结构的固体究竟是金属、半导体，还是实质上的绝缘体，我们回顾一下（§1.6）：一个布里渊区所包含的“允许的 k 矢量”——即彼此不同的电子波函数——恰与晶体中晶胞的数目一样多。而每个被允许的空间波函数可以容纳 2 个自旋相反的电子。于是，若单位体积中有 $N$ 个晶胞，则每条能带有 $2N$ 个电子态，即每条能带在每个晶胞摊到 2 个电子态。数一数每个晶胞中的电子数，我们就能对固体可能具有的性质作出若干估计。

(i) 每个晶胞有\emph{一个}自由电子的固体总是金属。典型的是一价金属——碱金属 Li、Na、K、Rb、Cs，以及贵金属（noble metals）Cu、Ag、Au。它们都恰好只有一条半满的能带，或者说半满的区。

(ii) 每个晶胞有\emph{奇数}个电子的固体总是金属。于是，在 Al、Ga、In、Tl 中每个原子有 3 个电子，只够填满一条半的能带。

但注意，As、Sb、Bi 每个原子有 5 个电子，它们结晶成的结构每个晶胞含 2 个原子，所以这条规则用不上。事实上，它们是半金属（semi-metal）：10 个电子几乎恰好填满 5 条能带，但第 5 条能带并未完全填满，还有少许\emph{交叠}进入第 6 条能带，因此不论温度如何，总有少量电子可以传导电流。

(iii) 每个晶胞有\emph{偶数}个电子的固体则不一定是绝缘体。这是因为各条能带在能量上可以交叠。让较低的能带留着不填满、而让一部分电子进入上面的能带，在能量上可能是有利的。

%FIG{68}: {图 68　相互交叠的能带。}

%FIG{69}: {图 69　二价金属的交叠能带：(a) 扩展区图式；(b) 简约区中能量随 k 的变化。}

在一维情形这不会发生：各条能带彼此完全分开（见图 38）。但在真实金属中，我们可能遇到图 69 所描绘的情况。虽然在任何特定地点越过区边界时都可能有能隙，但第 2 条能带的最低部分（例如在 $A$ 处）可能位于第 1 条能带最高部分之下，于是电子在 $A$ 处“溢出”区的边界，而区的角隅则空着。

这一点在扩展区图式的近自由电子（N.F.E.）模型中很容易看清（参见 §3.3）。一个容纳 $2N$ 个态的球面会在若干处与区边界相交。如果势（或赝势）相应的 Fourier 分量很小，其效果只不过是：当球面的各部分穿过每条区边界时把它们截断
%FIG{70}: {图 70　被区边界截断的自由电子球。}
——而费米面的其余部分实际上保持不变。对于通常的电学性质而言，费米面的这些碎片可以在简约区和重复区中重新连接起来，这一事实无关紧要。

(iv) 结果发现，所有\emph{二价}元素都是金属：能隙还不够大，不足以把电子都容纳在单一区内。但其中一些（Sr、Ba）是劣导体，其交叠可能很小。

(v) \emph{四价}元素或者是金属，或者是半导体。有趣的情形是 Sn，它在一种相中是金属，在另一种相中是半导体。晶体结构的改变会改变区的形状，从而改变出现足够大的能隙以容纳全部电子的可能性。在周期表第 IV 族中有一段有趣的递变：从 C——作为金刚石，它形成的半导体能隙如此之大，以致在实践中成了绝缘体——到 Si 和 Ge 这两种半导体，到既为金属又为半导体的 Sn，再到金属 Pb。所有这些元素的能带结构都可以用 N.F.E. 模型表示；半导体就是能隙很大且没有交叠的“金属”。

%FIG{71}: {图 71　(a) 过渡金属。(b) 贵金属。}

(vi) \emph{过渡元素}（transition elements），例如铁族的 Cr、Mn、Fe、Co、Ni，以及周期表中位置更高的其他各族，其特征在于内层 $d$ 壳层未满。例如在原子 Fe 中，尽管外层 $4s$ 态里还有另外 2 个电子，$3d$ 壳层的 10 个态中只有 6 个被填充。当原子聚到一起构成固体或液体时，这些外层价电子进入一条宽的“$s$ 能带”，与普通金属中的 N.F.E. 导带没有太大差别。在某些方面（见 §10.5），$d$ 电子的表现仍好像定域在各自的原子实中：对稀土族（rare earth group）的 $4f$ 电子来说，这实际上是一个极好的近似。但 $d$ 轨道实际上必须在原子之间相互交叠，以形成一条窄的“$d$ 能带”，每个原子至多可容纳 10 个电子。正如我们在 §3.10 中看到的，这条能带与 $s$ 能带交叉并杂化；严格说来，这种复杂的能带结构是与原先原子 $d$ 能级处的共振相联系的，不能分解成分开的 $s$ 能带与 $d$ 能带。

尽管如此，比如说，下述说法仍然是自然的：两条能带都不满，故材料呈金属性，导电主要依靠“$s$ 电子”（图 71(a)）。在贵金属中，$d$ 态实际上是满的，但它们在价电子的普通 N.F.E. 能带之内、费米能级以下几电子伏处生成一条共振 $d$ 能带（图 31(b)）。这对那些本会是简单一价金属的材料有显著影响（见 §9.4）。

\section{固体的类型：键图像}

用能带模型，我们所能轻松走到的程度也就到此为止了。但让我们再来看一看半导体 Si 和 Ge。它们具有“金刚石结构”（diamond structure），与 F.C.C.（面心立方）相像（见 §1.3），只不过
%FIG{72}: {图 72　金刚石结构中的四面体键。}
每个晶胞有两个原子。取一个以 $A$ 为角隅的 F.C.C. 格子，在沿立方体主对角线 $1/4$ 处的 $B$ 点放上另一个原子。让这个新原子作为第二个 F.C.C. 格子的角隅，与第一个相互贯穿。

这种结构与 F.C.C. 格子有相同的布里渊区。现在每个晶胞有 8 个电子，所以需要填满 4 条能带。如果从紧束缚模型的观点来看这 4 条能带，就会发现它们来自自由原子的 4 个轨道——Si 中的 $3s$ 轨道和三个 $3p$ 轨道；Ge 中的 $4s$ 与 $4p$ 轨道。当原子聚到一起时，这些轨道交叉并组合，形成 4 条能量范围大致相同的能带（我们在 §3.11 中注意到，若不是自旋–轨道相互作用，其中三条在区中心应当是简并的）。再上面一条能带与这些\emph{价带}（valence bands）之间隔着约 1 eV 量级的能隙。

但这种 $s$ 轨道与 $p$ 轨道的组合在化学键理论中是众所周知的。人们知道，这些波函数可以组合起来，给出四个轨道波函数，每一个都指向一个四面体的顶点。人们还知道，诸如 CH$_4$ 之类分子的四面体对称性与\emph{共价键}（covalent bond）的形成相联系，即通过在这些 $s$-$p$ 杂化轨道中共享电子对而成键。再看图 72，我们看到 $B$ 处的原子恰好有 4 个最近邻，排在四面体的顶点上。它们每一个又同样与其近邻相连，如此等等：换言之，\emph{整个晶体就像一个巨大的共价键合的分子}。

能带结构计算证实了这一点。用少数几个 O.P.W.（正交化平面波）配上合适的赝势系数，就能与实验能带结构符合得相当好。所得的波函数倾向于把电子密度集中到“键”上。我们示意地把电子结构表示成图 73 的样子。每个 Ge 或 Si 离子与近邻共享 8 个电子。这些电子其实并不定域在键上。像金属中的自由电子一样，它们的波函数扩展到整个晶体。但在相邻离子之间的区域中，总有一团高度集中的电子云，等效于一对自旋相反的电子。

%FIG{73}: {图 73　Ge 的电子结构（示意图）。}

然而要注意，这在某种程度上是一个自我实现的预言。在配位数低的格子中，muffin-tin 近似（§3.7）并不适用，电子的势能在间隙区域内并非处处均匀。在任何自洽计算的尝试（参见 §5.2）中，电子电荷都会自动沿相邻原子球之间的连线聚集，那里正是间隙势最低之处。如果我们一开始就要解释这类晶体结构何以出现，那么一般的化学论证看来是必不可少的。

现在有趣的是把这幅图像稍加改动：把 Ge 换成 InSb。我们记得 In 和 Sb 与 Ge 处在周期表的同一横行，所以它们具有相同的闭壳层基础芯。但 In 只有 3 个外层电子，Sb 有 5 个。它同样形成一种类金刚石
%FIG{74}: {图 74　锑化铟（示意图）。(a) 裸离子。(b) 带有对称的键。(c) 被离子上残余电荷极化的键电子。}
的晶格，只不过现在 In 占据全部 $A$ 位，而 Sb 占据全部 $B$ 位。示意地，我们得到如图 74(a) 那样的离子排列。

在许多方面它与 Ge 结构如此相像，以至于我们可以把每原子 8 个电子放入同样的“键-带”轨道，如图 74(b)。但这会留下一些残余电荷：Sb 上的电荷没有被它周围平均 4 个电子完全中和，而 In 上则多出一点额外的负电荷。不过，这些效应很容易通过让电荷云向 Sb 略微收缩、并离开 In 来容纳，如图 74(c)。这种化合物是半导体，基本能带结构与 Ge 相仿，尽管结果表明其能隙较小。它是典型的 \emph{III–V 族金属间化合物}（intermetallic compound）。

现在转到 ZnS，它是一种 \emph{II–VI 族化合物}——Zn 有 2 个电子，S 有 6 个。它形成的晶体与 InSb 完全相似——事实上，这就叫作\emph{闪锌矿}（zinc-blende）结构——而且也是半导体。但现在，电子对 $\mathrm{S}^{6+}$ 离子的静电吸引又因\emph{电子亲和势}（electron affinity）而增强：这个离子所能看到的 8 个电子倾向于收得更紧，力图构成一个完整的闭壳层，而让 $\mathrm{Zn}^{2+}$ 离子几乎不含电子。

%FIG{75}: {图 75　硫化锌（示意图）。}

这一过程在比如 CuCl 中走到了彻底的地步。Cu 原子完全失去它的电子，这个电子拿去构成 $\mathrm{Cl}^{-}$ 闭壳层。于是我们得到实际上相互独立的离子 $\mathrm{Cu}^{+}$ 和 $\mathrm{Cl}^{-}$，没有任何自由电子剩下。它将结晶成如图 76 那样的正负交替的晶格。但此时的内聚已不再与键相联系。它本质上来自离子间的静电吸引，可以按 §2.3 的方法计算。最近邻的四面体排列现在无关紧要了——要点是围绕一个负离子尽可能多地挤入正离子，反之亦然。\emph{离子型固体}倾向于结晶成 NaCl 和 CsCl 那样的结构。

对于这种晶体中的价电子，谈论“能带结构”究竟还有多大意义，是一个微妙的问题。这些电子的波函数紧紧定域在各自的离子上，不表现出诸如自由载流子“空穴”迁移率（参见 §6.6）那样的“Bloch 函数”性质；“能隙”大致就相当于金属原子的第二电离势。当然，被激发到这个隙之上的电子，其行为会更像普通导带中可迁移的载流子。

实际上，离子型固体与共价型固体之间并没有截然的分界。在 III–V 或 II–VI 族化合物中总存在某种程度的离子性。即便是离子晶体，也能表现出与半导体相似的性质。

%FIG{76}: {图 76　(a) 氯化铜（示意图）。(b) NaCl 结构。(c) CsCl 结构。}

这并没有穷尽固体的分类。还有\emph{分子型晶体}（molecular crystals），其构造单元是稳定的中性分子，例如 CH$_4$，靠彼此间的 van der Waals 引力凝聚在一起。还有一类特殊的\emph{氢键型结构}（hydrogen-bonded structures），其中分子基团之间的力通过一个共享的质子来传递。这些固体通常由分子的化学性质所支配。除了可作为简单晶体范式的稀有气体固体之外，它们往往结构复杂，其性质超出了本书的范围。
